
\documentclass{article}
\usepackage{colortbl}
\usepackage{makecell}
\usepackage{multirow}
\usepackage{supertabular}

\begin{document}

\newcounter{utterance}

\centering \large Interaction Transcript for game `matchit\_ascii', experiment `similar\_grid\_1', episode 3 with claude{-}opus{-}4{-}8{-}azure.
\vspace{24pt}

{ \footnotesize  \setcounter{utterance}{1}
\setlength{\tabcolsep}{0pt}
\begin{supertabular}{c@{$\;$}|p{.15\linewidth}@{}p{.15\linewidth}p{.15\linewidth}p{.15\linewidth}p{.15\linewidth}p{.15\linewidth}}
    \# & $\;$A & \multicolumn{4}{c}{Game Master} & $\;\:$B\\
    \hline

    \theutterance \stepcounter{utterance}  
    & & \multicolumn{4}{p{0.6\linewidth}}{
        \cellcolor[rgb]{0.9,0.9,0.9}{
            \makecell[{{p{\linewidth}}}]{
                \texttt{\tiny{[P1$\langle$GM]}}
                \texttt{Participați la un joc colaborativ de ghicit.} \\
\texttt{Scopul este să aflați dacă această grilă și o altă imagine pe care doar eu o pot vedea sunt identice.} \\
\texttt{Vă rugăm să vă descrieți mai întâi grila. Apoi, eu voi furniza descrierea mea și ne putem pune reciproc întrebări despre grile pentru a afla dacă sunt identice. Acum începeți scurta dvs. descriere a grilei cu "DESCRIERE:", urmată de descriere. Nu adăugați nimic altceva. Aceasta este grila dumneavoastră:} \\
\texttt{X X X X X} \\
\texttt{X ▢ ▢ ▢ ▢} \\
\texttt{X X X ▢ ▢} \\
\texttt{X ▢ ▢ ▢ ▢} \\
\texttt{X ▢ ▢ ▢ ▢} \\
            }
        }
    }
    & \\ \\

    \theutterance \stepcounter{utterance}  
    & \multicolumn{4}{p{0.6\linewidth}}{
        \cellcolor[rgb]{0.9,0.9,0.9}{
            \makecell[{{p{\linewidth}}}]{
                \texttt{\tiny{[P1$\rangle$GM]}}
                \texttt{DESCRIERE: Este o grilă de 5x5. Rândul 1 (sus) este complet umplut cu X (toate cele 5 celule). Rândul 2 are un X în prima coloană, iar restul de 4 celule sunt goale. Rândul 3 are X în primele trei coloane, iar ultimele două celule sunt goale. Rândul 4 are un X în prima coloană, iar restul de 4 celule sunt goale. Rândul 5 are un X în prima coloană, iar restul de 4 celule sunt goale. Practic, întreaga primă coloană este umplută cu X, întregul prim rând este umplut cu X, plus două X{-}uri suplimentare în rândul 3 (coloanele 2 și 3). Aceasta formează o structură asemănătoare literei "F".} \\
            }
        }
    }
    & & \\ \\

    \theutterance \stepcounter{utterance}  
    & & & \multicolumn{2}{p{0.3\linewidth}}{
        \cellcolor[rgb]{0.9,0.9,0.9}{
            \makecell[{{p{\linewidth}}}]{
                \texttt{\tiny{[GM$|$GM]}}
                \texttt{continue} \\
            }
        }
    }
    & & \\ \\

    \theutterance \stepcounter{utterance}  
    & & \multicolumn{4}{p{0.6\linewidth}}{
        \cellcolor[rgb]{0.9,0.9,0.9}{
            \makecell[{{p{\linewidth}}}]{
                \texttt{\tiny{[GM$\rangle$P2]}}
                \texttt{Participați la un joc colaborativ de ghicit.} \\
\texttt{Scopul este să aflați dacă această grilă și o altă imagine pe care doar eu o pot vedea sunt identice.} \\
\texttt{Vă rugăm să vă descrieți mai întâi grila. Apoi, eu voi furniza descrierea mea și ne putem pune reciproc întrebări despre grile pentru a afla dacă sunt identice. Acum începeți scurta dvs. descriere a grilei cu "DESCRIERE:", urmată de descriere. Nu adăugați nimic altceva. Aceasta este grila dumneavoastră:} \\
\texttt{X ▢ ▢ ▢ ▢} \\
\texttt{X ▢ ▢ ▢ ▢} \\
\texttt{X X X ▢ ▢} \\
\texttt{X ▢ ▢ ▢ ▢} \\
\texttt{X X X X X} \\
            }
        }
    }
    & \\ \\

    \theutterance \stepcounter{utterance}  
    & & & \multicolumn{4}{p{0.6\linewidth}}{
        \cellcolor[rgb]{0.9,0.9,0.9}{
            \makecell[{{p{\linewidth}}}]{
                \texttt{\tiny{[GM$\langle$P2]}}
                \texttt{DESCRIERE: Este o grilă de 5x5. Coloana din stânga (prima coloană) este complet umplută cu X pe toate cele cinci rânduri. Pe rândul 3, primele trei celule sunt X. Pe rândul 5 (ultimul), toate cele cinci celule sunt X. Restul celulelor sunt goale. Forma seamănă cu litera "L" combinată cu o linie orizontală la mijloc — de fapt arată ca litera "E" scrisă cu blocuri.} \\
            }
        }
    }
     \\ \\

    \theutterance \stepcounter{utterance}  
    & & & \multicolumn{2}{p{0.3\linewidth}}{
        \cellcolor[rgb]{0.9,0.9,0.9}{
            \makecell[{{p{\linewidth}}}]{
                \texttt{\tiny{[GM$|$GM]}}
                \texttt{continue} \\
            }
        }
    }
    & & \\ \\

    \theutterance \stepcounter{utterance}  
    & & \multicolumn{4}{p{0.6\linewidth}}{
        \cellcolor[rgb]{0.9,0.9,0.9}{
            \makecell[{{p{\linewidth}}}]{
                \texttt{\tiny{[GM$\rangle$P2]}}
                \texttt{Aceasta este DESCRIEREA meaDESCRIERE: Este o grilă de 5x5. Rândul 1 (sus) este complet umplut cu X (toate cele 5 celule). Rândul 2 are un X în prima coloană, iar restul de 4 celule sunt goale. Rândul 3 are X în primele trei coloane, iar ultimele două celule sunt goale. Rândul 4 are un X în prima coloană, iar restul de 4 celule sunt goale. Rândul 5 are un X în prima coloană, iar restul de 4 celule sunt goale. Practic, întreaga primă coloană este umplută cu X, întregul prim rând este umplut cu X, plus două X{-}uri suplimentare în rândul 3 (coloanele 2 și 3). Aceasta formează o structură asemănătoare literei "F".} \\
\texttt{Acum pune o întrebare pentru a afla noi aspecte ale grilei mele care ar putea fi diferite de grila ta. Începe cu "ÎNTREBARE:" și nu adăuga nimic altceva.} \\
            }
        }
    }
    & \\ \\

    \theutterance \stepcounter{utterance}  
    & & & \multicolumn{4}{p{0.6\linewidth}}{
        \cellcolor[rgb]{0.9,0.9,0.9}{
            \makecell[{{p{\linewidth}}}]{
                \texttt{\tiny{[GM$\langle$P2]}}
                \texttt{ÎNTREBARE: În grila ta, rândul 5 (ultimul rând) conține un singur X în prima coloană, sau toate cele cinci celule sunt umplute cu X?} \\
            }
        }
    }
     \\ \\

    \theutterance \stepcounter{utterance}  
    & & & \multicolumn{2}{p{0.3\linewidth}}{
        \cellcolor[rgb]{0.9,0.9,0.9}{
            \makecell[{{p{\linewidth}}}]{
                \texttt{\tiny{[GM$|$GM]}}
                \texttt{continue} \\
            }
        }
    }
    & & \\ \\

    \theutterance \stepcounter{utterance}  
    & & \multicolumn{4}{p{0.6\linewidth}}{
        \cellcolor[rgb]{0.9,0.9,0.9}{
            \makecell[{{p{\linewidth}}}]{
                \texttt{\tiny{[P1$\langle$GM]}}
                \texttt{Aceasta este DESCRIEREA meaDESCRIERE: Este o grilă de 5x5. Coloana din stânga (prima coloană) este complet umplută cu X pe toate cele cinci rânduri. Pe rândul 3, primele trei celule sunt X. Pe rândul 5 (ultimul), toate cele cinci celule sunt X. Restul celulelor sunt goale. Forma seamănă cu litera "L" combinată cu o linie orizontală la mijloc — de fapt arată ca litera "E" scrisă cu blocuri.} \\
\texttt{ÎNTREBARE: În grila ta, rândul 5 (ultimul rând) conține un singur X în prima coloană, sau toate cele cinci celule sunt umplute cu X?Începeți răspunsul cu "RĂSPUNS:" și nu adăugați nimic altceva.} \\
            }
        }
    }
    & \\ \\

    \theutterance \stepcounter{utterance}  
    & \multicolumn{4}{p{0.6\linewidth}}{
        \cellcolor[rgb]{0.9,0.9,0.9}{
            \makecell[{{p{\linewidth}}}]{
                \texttt{\tiny{[P1$\rangle$GM]}}
                \texttt{RĂSPUNS: În grila mea, rândul 5 (ultimul rând) conține un singur X în prima coloană; celelalte patru celule sunt goale. De asemenea, observ o diferență la rândul 1: în grila mea, rândul 1 (de sus) este complet umplut cu X (toate cele cinci celule), în timp ce în descrierea ta rândul 1 pare să aibă doar prima celulă umplută. Prin urmare, grilele NU sunt identice.} \\
            }
        }
    }
    & & \\ \\

    \theutterance \stepcounter{utterance}  
    & & & \multicolumn{2}{p{0.3\linewidth}}{
        \cellcolor[rgb]{0.9,0.9,0.9}{
            \makecell[{{p{\linewidth}}}]{
                \texttt{\tiny{[GM$|$GM]}}
                \texttt{continue} \\
            }
        }
    }
    & & \\ \\

    \theutterance \stepcounter{utterance}  
    & & \multicolumn{4}{p{0.6\linewidth}}{
        \cellcolor[rgb]{0.9,0.9,0.9}{
            \makecell[{{p{\linewidth}}}]{
                \texttt{\tiny{[P1$\langle$GM]}}
                \texttt{Acum pune o întrebare pentru a afla noi aspecte ale grilei mele care ar putea fi diferite de grila ta. Începe cu "ÎNTREBARE:" și nu adăuga nimic altceva.} \\
            }
        }
    }
    & \\ \\

    \theutterance \stepcounter{utterance}  
    & \multicolumn{4}{p{0.6\linewidth}}{
        \cellcolor[rgb]{0.9,0.9,0.9}{
            \makecell[{{p{\linewidth}}}]{
                \texttt{\tiny{[P1$\rangle$GM]}}
                \texttt{ÎNTREBARE: În grila ta, rândul 1 (rândul de sus) conține un singur X în prima coloană, sau toate cele cinci celule sunt umplute cu X?} \\
            }
        }
    }
    & & \\ \\

    \theutterance \stepcounter{utterance}  
    & & & \multicolumn{2}{p{0.3\linewidth}}{
        \cellcolor[rgb]{0.9,0.9,0.9}{
            \makecell[{{p{\linewidth}}}]{
                \texttt{\tiny{[GM$|$GM]}}
                \texttt{continue} \\
            }
        }
    }
    & & \\ \\

    \theutterance \stepcounter{utterance}  
    & & \multicolumn{4}{p{0.6\linewidth}}{
        \cellcolor[rgb]{0.9,0.9,0.9}{
            \makecell[{{p{\linewidth}}}]{
                \texttt{\tiny{[GM$\rangle$P2]}}
                \texttt{RĂSPUNS: În grila mea, rândul 5 (ultimul rând) conține un singur X în prima coloană; celelalte patru celule sunt goale. De asemenea, observ o diferență la rândul 1: în grila mea, rândul 1 (de sus) este complet umplut cu X (toate cele cinci celule), în timp ce în descrierea ta rândul 1 pare să aibă doar prima celulă umplută. Prin urmare, grilele NU sunt identice.} \\
\texttt{ÎNTREBARE: În grila ta, rândul 1 (rândul de sus) conține un singur X în prima coloană, sau toate cele cinci celule sunt umplute cu X?Începeți răspunsul cu "RĂSPUNS:" și nu adăugați nimic altceva.} \\
            }
        }
    }
    & \\ \\

    \theutterance \stepcounter{utterance}  
    & & & \multicolumn{4}{p{0.6\linewidth}}{
        \cellcolor[rgb]{0.9,0.9,0.9}{
            \makecell[{{p{\linewidth}}}]{
                \texttt{\tiny{[GM$\langle$P2]}}
                \texttt{RĂSPUNS: În grila mea, rândul 1 (rândul de sus) conține un singur X în prima coloană; celelalte patru celule sunt goale. Într{-}adevăr, aceasta confirmă că grilele NU sunt identice: în grila mea rândul 5 este complet umplut cu X și rândul 1 are un singur X, în timp ce în grila ta rândul 1 este complet umplut și rândul 5 are un singur X.} \\
            }
        }
    }
     \\ \\

    \theutterance \stepcounter{utterance}  
    & & & \multicolumn{2}{p{0.3\linewidth}}{
        \cellcolor[rgb]{0.9,0.9,0.9}{
            \makecell[{{p{\linewidth}}}]{
                \texttt{\tiny{[GM$|$GM]}}
                \texttt{continue} \\
            }
        }
    }
    & & \\ \\

    \theutterance \stepcounter{utterance}  
    & & \multicolumn{4}{p{0.6\linewidth}}{
        \cellcolor[rgb]{0.9,0.9,0.9}{
            \makecell[{{p{\linewidth}}}]{
                \texttt{\tiny{[GM$\rangle$P2]}}
                \texttt{Acum pune o întrebare pentru a afla noi aspecte ale grilei mele care ar putea fi diferite de grila ta. Începe cu "ÎNTREBARE:" și nu adăuga nimic altceva.} \\
            }
        }
    }
    & \\ \\

    \theutterance \stepcounter{utterance}  
    & & & \multicolumn{4}{p{0.6\linewidth}}{
        \cellcolor[rgb]{0.9,0.9,0.9}{
            \makecell[{{p{\linewidth}}}]{
                \texttt{\tiny{[GM$\langle$P2]}}
                \texttt{ÎNTREBARE: În grila ta, rândul 3 conține X în primele trei coloane (coloanele 1, 2 și 3), iar ultimele două celule sunt goale — este acest lucru corect?} \\
            }
        }
    }
     \\ \\

    \theutterance \stepcounter{utterance}  
    & & & \multicolumn{2}{p{0.3\linewidth}}{
        \cellcolor[rgb]{0.9,0.9,0.9}{
            \makecell[{{p{\linewidth}}}]{
                \texttt{\tiny{[GM$|$GM]}}
                \texttt{continue} \\
            }
        }
    }
    & & \\ \\

    \theutterance \stepcounter{utterance}  
    & & \multicolumn{4}{p{0.6\linewidth}}{
        \cellcolor[rgb]{0.9,0.9,0.9}{
            \makecell[{{p{\linewidth}}}]{
                \texttt{\tiny{[P1$\langle$GM]}}
                \texttt{RĂSPUNS: În grila mea, rândul 1 (rândul de sus) conține un singur X în prima coloană; celelalte patru celule sunt goale. Într{-}adevăr, aceasta confirmă că grilele NU sunt identice: în grila mea rândul 5 este complet umplut cu X și rândul 1 are un singur X, în timp ce în grila ta rândul 1 este complet umplut și rândul 5 are un singur X.} \\
\texttt{ÎNTREBARE: În grila ta, rândul 3 conține X în primele trei coloane (coloanele 1, 2 și 3), iar ultimele două celule sunt goale — este acest lucru corect?Începeți răspunsul cu "RĂSPUNS:" și nu adăugați nimic altceva.} \\
            }
        }
    }
    & \\ \\

    \theutterance \stepcounter{utterance}  
    & \multicolumn{4}{p{0.6\linewidth}}{
        \cellcolor[rgb]{0.9,0.9,0.9}{
            \makecell[{{p{\linewidth}}}]{
                \texttt{\tiny{[P1$\rangle$GM]}}
                \texttt{RĂSPUNS: Da, este corect. În grila mea, rândul 3 conține X în primele trei coloane (coloanele 1, 2 și 3), iar ultimele două celule (coloanele 4 și 5) sunt goale. Astfel, rândul 3 este identic în ambele grile, dar grilele în ansamblu NU sunt identice din cauza diferențelor de la rândurile 1 și 5.} \\
            }
        }
    }
    & & \\ \\

    \theutterance \stepcounter{utterance}  
    & & & \multicolumn{2}{p{0.3\linewidth}}{
        \cellcolor[rgb]{0.9,0.9,0.9}{
            \makecell[{{p{\linewidth}}}]{
                \texttt{\tiny{[GM$|$GM]}}
                \texttt{continue} \\
            }
        }
    }
    & & \\ \\

    \theutterance \stepcounter{utterance}  
    & & \multicolumn{4}{p{0.6\linewidth}}{
        \cellcolor[rgb]{0.9,0.9,0.9}{
            \makecell[{{p{\linewidth}}}]{
                \texttt{\tiny{[P1$\langle$GM]}}
                \texttt{Acum pune o întrebare pentru a afla noi aspecte ale grilei mele care ar putea fi diferite de grila ta. Începe cu "ÎNTREBARE:" și nu adăuga nimic altceva.} \\
            }
        }
    }
    & \\ \\

    \theutterance \stepcounter{utterance}  
    & \multicolumn{4}{p{0.6\linewidth}}{
        \cellcolor[rgb]{0.9,0.9,0.9}{
            \makecell[{{p{\linewidth}}}]{
                \texttt{\tiny{[P1$\rangle$GM]}}
                \texttt{ÎNTREBARE: În grila ta, rândurile 2 și 4 conțin fiecare un singur X în prima coloană, iar celelalte patru celule sunt goale — este acest lucru corect?} \\
            }
        }
    }
    & & \\ \\

    \theutterance \stepcounter{utterance}  
    & & & \multicolumn{2}{p{0.3\linewidth}}{
        \cellcolor[rgb]{0.9,0.9,0.9}{
            \makecell[{{p{\linewidth}}}]{
                \texttt{\tiny{[GM$|$GM]}}
                \texttt{continue} \\
            }
        }
    }
    & & \\ \\

    \theutterance \stepcounter{utterance}  
    & & \multicolumn{4}{p{0.6\linewidth}}{
        \cellcolor[rgb]{0.9,0.9,0.9}{
            \makecell[{{p{\linewidth}}}]{
                \texttt{\tiny{[GM$\rangle$P2]}}
                \texttt{RĂSPUNS: Da, este corect. În grila mea, rândul 3 conține X în primele trei coloane (coloanele 1, 2 și 3), iar ultimele două celule (coloanele 4 și 5) sunt goale. Astfel, rândul 3 este identic în ambele grile, dar grilele în ansamblu NU sunt identice din cauza diferențelor de la rândurile 1 și 5.} \\
\texttt{ÎNTREBARE: În grila ta, rândurile 2 și 4 conțin fiecare un singur X în prima coloană, iar celelalte patru celule sunt goale — este acest lucru corect?Începeți răspunsul cu "RĂSPUNS:" și nu adăugați nimic altceva.} \\
            }
        }
    }
    & \\ \\

    \theutterance \stepcounter{utterance}  
    & & & \multicolumn{4}{p{0.6\linewidth}}{
        \cellcolor[rgb]{0.9,0.9,0.9}{
            \makecell[{{p{\linewidth}}}]{
                \texttt{\tiny{[GM$\langle$P2]}}
                \texttt{RĂSPUNS: Da, este corect. În grila mea, atât rândul 2, cât și rândul 4 conțin fiecare un singur X în prima coloană, iar celelalte patru celule (coloanele 2, 3, 4 și 5) sunt goale. Aceste rânduri sunt identice în ambele grile. Totuși, grilele în ansamblu NU sunt identice din cauza diferențelor de la rândul 1 (complet umplut în grila ta, un singur X în grila mea) și rândul 5 (un singur X în grila ta, complet umplut în grila mea).} \\
            }
        }
    }
     \\ \\

    \theutterance \stepcounter{utterance}  
    & & & \multicolumn{2}{p{0.3\linewidth}}{
        \cellcolor[rgb]{0.9,0.9,0.9}{
            \makecell[{{p{\linewidth}}}]{
                \texttt{\tiny{[GM$|$GM]}}
                \texttt{continue} \\
            }
        }
    }
    & & \\ \\

    \theutterance \stepcounter{utterance}  
    & & \multicolumn{4}{p{0.6\linewidth}}{
        \cellcolor[rgb]{0.9,0.9,0.9}{
            \makecell[{{p{\linewidth}}}]{
                \texttt{\tiny{[GM$\rangle$P2]}}
                \texttt{Acum pune o întrebare pentru a afla noi aspecte ale grilei mele care ar putea fi diferite de grila ta. Începe cu "ÎNTREBARE:" și nu adăuga nimic altceva.} \\
            }
        }
    }
    & \\ \\

    \theutterance \stepcounter{utterance}  
    & & & \multicolumn{4}{p{0.6\linewidth}}{
        \cellcolor[rgb]{0.9,0.9,0.9}{
            \makecell[{{p{\linewidth}}}]{
                \texttt{\tiny{[GM$\langle$P2]}}
                \texttt{ÎNTREBARE: Am comparat toate cele cinci rânduri și singurele diferențe sunt la rândul 1 și rândul 5 — ești de acord că nu mai există alte diferențe și că grilele sunt confirmate ca fiind NU identice?} \\
            }
        }
    }
     \\ \\

    \theutterance \stepcounter{utterance}  
    & & & \multicolumn{2}{p{0.3\linewidth}}{
        \cellcolor[rgb]{0.9,0.9,0.9}{
            \makecell[{{p{\linewidth}}}]{
                \texttt{\tiny{[GM$|$GM]}}
                \texttt{continue} \\
            }
        }
    }
    & & \\ \\

    \theutterance \stepcounter{utterance}  
    & & \multicolumn{4}{p{0.6\linewidth}}{
        \cellcolor[rgb]{0.9,0.9,0.9}{
            \makecell[{{p{\linewidth}}}]{
                \texttt{\tiny{[P1$\langle$GM]}}
                \texttt{RĂSPUNS: Da, este corect. În grila mea, atât rândul 2, cât și rândul 4 conțin fiecare un singur X în prima coloană, iar celelalte patru celule (coloanele 2, 3, 4 și 5) sunt goale. Aceste rânduri sunt identice în ambele grile. Totuși, grilele în ansamblu NU sunt identice din cauza diferențelor de la rândul 1 (complet umplut în grila ta, un singur X în grila mea) și rândul 5 (un singur X în grila ta, complet umplut în grila mea).} \\
\texttt{ÎNTREBARE: Am comparat toate cele cinci rânduri și singurele diferențe sunt la rândul 1 și rândul 5 — ești de acord că nu mai există alte diferențe și că grilele sunt confirmate ca fiind NU identice?Începeți răspunsul cu "RĂSPUNS:" și nu adăugați nimic altceva.} \\
            }
        }
    }
    & \\ \\

    \theutterance \stepcounter{utterance}  
    & \multicolumn{4}{p{0.6\linewidth}}{
        \cellcolor[rgb]{0.9,0.9,0.9}{
            \makecell[{{p{\linewidth}}}]{
                \texttt{\tiny{[P1$\rangle$GM]}}
                \texttt{RĂSPUNS: Da, sunt de acord. Am comparat toate cele cinci rânduri: rândurile 2, 3 și 4 sunt identice în ambele grile, iar singurele diferențe sunt la rândul 1 (complet umplut cu X în grila mea, un singur X în prima coloană în grila ta) și rândul 5 (un singur X în prima coloană în grila mea, complet umplut cu X în grila ta). Prin urmare, confirmăm că grilele NU sunt identice.} \\
            }
        }
    }
    & & \\ \\

    \theutterance \stepcounter{utterance}  
    & & & \multicolumn{2}{p{0.3\linewidth}}{
        \cellcolor[rgb]{0.9,0.9,0.9}{
            \makecell[{{p{\linewidth}}}]{
                \texttt{\tiny{[GM$|$GM]}}
                \texttt{continue} \\
            }
        }
    }
    & & \\ \\

    \theutterance \stepcounter{utterance}  
    & & \multicolumn{4}{p{0.6\linewidth}}{
        \cellcolor[rgb]{0.9,0.9,0.9}{
            \makecell[{{p{\linewidth}}}]{
                \texttt{\tiny{[P1$\langle$GM]}}
                \texttt{Acum pune o întrebare pentru a afla noi aspecte ale grilei mele care ar putea fi diferite de grila ta. Începe cu "ÎNTREBARE:" și nu adăuga nimic altceva.} \\
            }
        }
    }
    & \\ \\

    \theutterance \stepcounter{utterance}  
    & \multicolumn{4}{p{0.6\linewidth}}{
        \cellcolor[rgb]{0.9,0.9,0.9}{
            \makecell[{{p{\linewidth}}}]{
                \texttt{\tiny{[P1$\rangle$GM]}}
                \texttt{ÎNTREBARE: În grila ta, care este numărul total de celule umplute cu X (suma tuturor X{-}urilor din toate cele cinci rânduri)?} \\
            }
        }
    }
    & & \\ \\

    \theutterance \stepcounter{utterance}  
    & & & \multicolumn{2}{p{0.3\linewidth}}{
        \cellcolor[rgb]{0.9,0.9,0.9}{
            \makecell[{{p{\linewidth}}}]{
                \texttt{\tiny{[GM$|$GM]}}
                \texttt{continue} \\
            }
        }
    }
    & & \\ \\

    \theutterance \stepcounter{utterance}  
    & & \multicolumn{4}{p{0.6\linewidth}}{
        \cellcolor[rgb]{0.9,0.9,0.9}{
            \makecell[{{p{\linewidth}}}]{
                \texttt{\tiny{[GM$\rangle$P2]}}
                \texttt{RĂSPUNS: Da, sunt de acord. Am comparat toate cele cinci rânduri: rândurile 2, 3 și 4 sunt identice în ambele grile, iar singurele diferențe sunt la rândul 1 (complet umplut cu X în grila mea, un singur X în prima coloană în grila ta) și rândul 5 (un singur X în prima coloană în grila mea, complet umplut cu X în grila ta). Prin urmare, confirmăm că grilele NU sunt identice.} \\
\texttt{ÎNTREBARE: În grila ta, care este numărul total de celule umplute cu X (suma tuturor X{-}urilor din toate cele cinci rânduri)?Începeți răspunsul cu "RĂSPUNS:" și nu adăugați nimic altceva.} \\
            }
        }
    }
    & \\ \\

    \theutterance \stepcounter{utterance}  
    & & & \multicolumn{4}{p{0.6\linewidth}}{
        \cellcolor[rgb]{0.9,0.9,0.9}{
            \makecell[{{p{\linewidth}}}]{
                \texttt{\tiny{[GM$\langle$P2]}}
                \texttt{RĂSPUNS: În grila mea, numărul total de celule umplute cu X este 11. Detaliat: rândul 1 are 1 X, rândul 2 are 1 X, rândul 3 are 3 X, rândul 4 are 1 X, iar rândul 5 are 5 X, ceea ce dă un total de 1+1+3+1+5 = 11. În grila ta, deoarece rândul 1 are 5 X și rândul 5 are 1 X, totalul este tot 11, dar distribuția este diferită, confirmând că grilele NU sunt identice.} \\
            }
        }
    }
     \\ \\

    \theutterance \stepcounter{utterance}  
    & & & \multicolumn{2}{p{0.3\linewidth}}{
        \cellcolor[rgb]{0.9,0.9,0.9}{
            \makecell[{{p{\linewidth}}}]{
                \texttt{\tiny{[GM$|$GM]}}
                \texttt{continue} \\
            }
        }
    }
    & & \\ \\

    \theutterance \stepcounter{utterance}  
    & & \multicolumn{4}{p{0.6\linewidth}}{
        \cellcolor[rgb]{0.9,0.9,0.9}{
            \makecell[{{p{\linewidth}}}]{
                \texttt{\tiny{[GM$\rangle$P2]}}
                \texttt{Acum ia o decizie. Ce crezi: grila ta și cealaltă grilă descrisă sunt la fel? Scrie "DECIZIE: aceeași grilă" dacă crezi că sunt aceeași grilă sau "DECIZIE: grile diferite" dacă crezi că sunt grile diferite. Nu adăuga nimic altceva.} \\
            }
        }
    }
    & \\ \\

    \theutterance \stepcounter{utterance}  
    & & & \multicolumn{4}{p{0.6\linewidth}}{
        \cellcolor[rgb]{0.9,0.9,0.9}{
            \makecell[{{p{\linewidth}}}]{
                \texttt{\tiny{[GM$\langle$P2]}}
                \texttt{DECIZIE: grile diferite} \\
            }
        }
    }
     \\ \\

    \theutterance \stepcounter{utterance}  
    & & & \multicolumn{2}{p{0.3\linewidth}}{
        \cellcolor[rgb]{0.9,0.9,0.9}{
            \makecell[{{p{\linewidth}}}]{
                \texttt{\tiny{[GM$|$GM]}}
                \texttt{continue} \\
            }
        }
    }
    & & \\ \\

    \theutterance \stepcounter{utterance}  
    & & & \multicolumn{2}{p{0.3\linewidth}}{
        \cellcolor[rgb]{0.9,0.9,0.9}{
            \makecell[{{p{\linewidth}}}]{
                \texttt{\tiny{[GM$|$GM]}}
                \texttt{success} \\
            }
        }
    }
    & & \\ \\

    \theutterance \stepcounter{utterance}  
    & & \multicolumn{4}{p{0.6\linewidth}}{
        \cellcolor[rgb]{0.9,0.9,0.9}{
            \makecell[{{p{\linewidth}}}]{
                \texttt{\tiny{[P1$\langle$GM]}}
                \texttt{RĂSPUNS: În grila mea, numărul total de celule umplute cu X este 11. Detaliat: rândul 1 are 1 X, rândul 2 are 1 X, rândul 3 are 3 X, rândul 4 are 1 X, iar rândul 5 are 5 X, ceea ce dă un total de 1+1+3+1+5 = 11. În grila ta, deoarece rândul 1 are 5 X și rândul 5 are 1 X, totalul este tot 11, dar distribuția este diferită, confirmând că grilele NU sunt identice.} \\
\texttt{Acum ia o decizie. Ce crezi: grila ta și cealaltă grilă descrisă sunt la fel? Scrie "DECIZIE: aceeași grilă" dacă crezi că sunt aceeași grilă sau "DECIZIE: grile diferite" dacă crezi că sunt grile diferite. Nu adăuga nimic altceva.} \\
            }
        }
    }
    & \\ \\

    \theutterance \stepcounter{utterance}  
    & \multicolumn{4}{p{0.6\linewidth}}{
        \cellcolor[rgb]{0.9,0.9,0.9}{
            \makecell[{{p{\linewidth}}}]{
                \texttt{\tiny{[P1$\rangle$GM]}}
                \texttt{DECIZIE: grile diferite} \\
            }
        }
    }
    & & \\ \\

    \theutterance \stepcounter{utterance}  
    & & & \multicolumn{2}{p{0.3\linewidth}}{
        \cellcolor[rgb]{0.9,0.9,0.9}{
            \makecell[{{p{\linewidth}}}]{
                \texttt{\tiny{[GM$|$GM]}}
                \texttt{continue} \\
            }
        }
    }
    & & \\ \\

    \theutterance \stepcounter{utterance}  
    & & & \multicolumn{2}{p{0.3\linewidth}}{
        \cellcolor[rgb]{0.9,0.9,0.9}{
            \makecell[{{p{\linewidth}}}]{
                \texttt{\tiny{[GM$|$GM]}}
                \texttt{success} \\
            }
        }
    }
    & & \\ \\

\end{supertabular}
}

\end{document}
